\section{Discussion}
\label{sec:discussion}

\subsection{RL vs.\ Rule-Based Handover Policies}
The RL-Optimised policy achieves the lowest median handover latency among
the three compared algorithms (see Table~\ref{tab:handover_kpi}), largely
because the DQN agent learns to anticipate signal degradation from the
RSRP/RSRQ trends rather than reacting only to threshold crossings.
The RSRP-Based policy performs better than the Distance-Based one in mobile
scenarios because received power is a more reliable proxy for link quality
than geometry alone.

\subsection{Federated Learning Benefits}
The FL xApp converges to a global prediction model without requiring UE
measurements to leave the gNB, in line with data-minimisation
requirements of O-RAN privacy guidelines.
The digital-twin gating further improves effective performance by filtering
RL actions that are predicted to degrade KPIs, acting as a safety
mechanism during early training stages.

\subsection{MEC Offloading}
AR/VR Streaming and Video Analytics exhibit the highest CPU utilisation,
consistent with their compute-intensive processing pipelines.
The Federated Learning service benefits disproportionately from high-memory
nodes, suggesting that MEC host selection should consider both CPU and
memory when placing FL aggregation tasks.

\subsection{Limitations}
\begin{itemize}
  \item The current THz channel model does not incorporate molecular
        absorption coefficients; this is a planned extension.
  \item RL training is performed within the simulation loop; a
        production deployment would require offline pre-training.
  \item Results are based on synthetic traffic models; real-world
        traces would strengthen the evaluation.
\end{itemize}

\subsection{Future Work}
% TODO: expand with specific planned experiments.
Future directions include: (i) multi-agent RL across gNBs, (ii) non-IID
FedAvg variants (FedProx), (iii) full THz molecular absorption modelling,
and (iv) integration with an open-source near-RT RIC (e.g., OpenAirInterface
FlexRIC).
